%% File: chapters/chapter4.tex
%% Author: Y.U.P. (yashparitkar)
%%
%% Description: Integrating the generated core into an FPGA project.

\documentclass[../manual.tex]{subfiles}

\begin{document}

\chapter{Integrating LibreILA into a Design}
\label{ch:integration}

\todo{What to do with the generated code once it exists.}

\section{Files to add}
\label{sec:integ_files}

\todo{What the generator emits for GEN\_TYPE 0 and 1, and which of those files go into the project. VHDL Style Guide can be run over them first where a project has a house style.}

\section{Core or wrapper}
\label{sec:integ_which}

\todo{Choosing between the two, and what each one costs and buys.}

\section{Splicing the probe into a link}
\label{sec:integ_splice}

\todo{What the probe master and probe slave sides are, and where the core sits in the link.}

\section{Clocks and reset}
\label{sec:integ_clocks}

\todo{\texttt{samp\_aclk} takes the clock the probed link runs on, not a convenient free-running one; \texttt{s\_axil\_aclk} takes whatever the bus runs on. The practical reset rule, with the requirements themselves left in the datasheet.}

\section{Timing constraints}
\label{sec:integ_sdc}

\todo{Work in progress: not verified in SDC yet.}

\section{Connecting the UART}
\label{sec:integ_uart}

\todo{Any two pins will do, chosen in the design constraints.}

\section{Checking the build}
\label{sec:integ_check}

\todo{Confirming the sample buffer was inferred as RAM and that timing is met.}

\end{document}
